\section{The Residual Stream as the Belief State}
\label{sec:residual-stream}

The starting point for seeing the Bayes net in the transformer is the residual
stream. Everything else follows from getting this right.

In the standard transformer encoder, each token position carries a vector of
dimension $D_{\mathrm{model}}$ that persists across all layers. Each layer reads
from this vector and writes its output back into it via a residual connection.
Nothing is discarded; the stream accumulates. Elhage et al.~\cite{elhage2021mathematical}
name this the \emph{residual stream} and frame it as a shared workspace: attention
heads and feed-forward layers are read-write operations on a common scratchpad.

In the Bayes net view, the residual stream at token position $t$ is the
\emph{belief state} of factor graph node $t$. Specifically, the formal token
encoding from~\cite{coppola2026transformers} assigns clean semantic roles to each
dimension:

\begin{center}
\begin{tabular}{cl}
\toprule
Dimensions & Content \\
\midrule
0 & own belief (initialized to $0.5$) \\
1--4 & factor table entries $[f_{00}, f_{01}, f_{10}, f_{11}]$ \\
5 & node type (0 = variable, 1 = factor) \\
6 & own index $/ (n-1)$ \\
7 & neighbor index $/ (n-1)$ \\
\bottomrule
\end{tabular}
\end{center}

Dimensions 5--7 form the \emph{routing key}: the complete inferential identity of
the token. Dimensions 0--4 are continuous parameters: they supply the magnitudes
the inference computes over, but they do not determine the structure of the
computation. This split is not a convenience of presentation --- it is a theorem.
Two tokens with the same routing key receive identical attention patterns
regardless of their continuous parameters~\cite{coppola2026transformers}.

\subsection*{Routing Key = Concept. Continuous Dimensions = Magnitude.}

This is the right level of abstraction for thinking about what the residual stream
represents. The routing key identifies \emph{which} concept a token corresponds to.
The continuous dimensions encode \emph{how strongly} the evidence supports that
concept. A production transformer with $D_{\mathrm{model}} = 4096$ is not tracking
4096 independent propositions --- it is tracking a much smaller number of concepts
(bounded by the routing key space) while using the remaining dimensions to encode
magnitudes and handle superposition.

The residual connection is what makes this interpretation coherent across layers.
Belief propagation is an iterative algorithm: each round refines the belief at every
node by incorporating messages from neighbors. The residual connection ensures that
each layer's update is added to the accumulated belief, not computed from scratch.
Layer $l$ does not replace the belief; it refines it.

\subsection*{What This Looks Like in Trained Models}

The mechanistic interpretability literature has independently converged on a
compatible picture. Elhage et al.~\cite{elhage2022superposition} show that the
residual stream represents far more features than $D_{\mathrm{model}}$ dimensions
via superposition: features are encoded as nearly-orthogonal directions in a
high-dimensional space, with interference managed by sparsity. In the Bayes net
view, this means the residual stream is simultaneously tracking many factor graph
nodes, each encoded as a direction rather than a dedicated dimension.

Sparse autoencoder (SAE) features~\cite{templeton2024monosemanticity} are the
natural candidates for individual factor graph nodes. A monosemantic SAE feature ---
one that activates for a single coherent concept --- corresponds to a single routing
key in the formal construction. The feature direction in residual stream space is
where the belief value for that concept lives.

The broader implication: the residual stream is not a flat representation. It has
structure --- routing structure, imposed by the attention weights, that determines
which features interact with which. That structure is the implicit factor graph.